\section{课程定位：评估既是测量科学，也是 Agent 基础设施}

这讲同时承担两项任务：建立 LLM Agent evaluation 的概念框架，并把课程项目具体化为可运行的 Green Agent。两部分不是行政内容与技术内容的拼接。一个 benchmark 只有被实现成可重复启动环境、驱动 participant agent、收集轨迹并计算指标的系统，才真正具备可用性；反过来，项目代码若没有 outcome validity、污染控制和偏差分析，也只是一个会返回数字的脚本。

\subsection{为什么 Agent 时代更需要评估}

slides 3--8 从最基本的问题开始。evaluation 是对模型和 Agent 的系统性、可重复测量，用来判断能力是否进步、比较方案并决定能否部署。它贯穿训练、开发、模型选择和上线监控。传统 LLM 多为静态 text-to-text；Agent 具有工具、记忆、多步推理和环境交互，输出不再只是文本，而是会改变外部状态的轨迹。因此一次成功截图无法证明系统可靠，必须复现任务、环境和评分过程。

\lecturefigure{slides-images/slide-003.png}{官方 slide 3：LLM Agent evaluation 为什么重要}
\lecturefigure{slides-images/slide-004.png}{官方 slide 4：evaluation 的系统性与可重复性定义}
\lecturefigure{slides-images/slide-005.png}{官方 slide 5：训练、开发、选择与部署阶段都需要 evaluation}
\lecturefigure{slides-images/slide-006.png}{官方 slide 6：公平比较、弱点发现与科学进步}
\lecturefigure{slides-images/slide-008.png}{官方 slide 8：从静态 LLM eval 到有工具/记忆/多步环境的 Agent eval}

\begin{importantbox}{Evaluation 是决策接口}
评估的产出不应只是 leaderboard 分数，而应支持四类决策：是否继续训练、选择哪个模型/Agent、哪个失败模式优先修复、是否达到部署门槛。每个指标都要绑定一个可执行动作；无法改变决策的数字只能算 telemetry，不能算完整 evaluation。
\end{importantbox}

\begin{knowledgebox}{读图：从 output measurement 到 trajectory measurement}
静态 LLM eval 通常固定输入并比较输出；Agent eval 还要记录初始状态、可用工具、每一步动作、工具返回、环境变化、终止条件和成本。若只对最终文字评分，模型可能在中间越权、破坏状态或依赖偶然缓存。slide 8 的本质是把评估单位从 response 扩展为 trajectory 与 outcome。
\end{knowledgebox}

\textbf{课堂提示：}授课团队在开场明确说明，本讲会直接连接 phase-one project。原因是 Green Agent 本身就是 evaluator/orchestrator：它要准备环境、测试 White Agent、验证结果并汇报。项目的工程边界因此来自评估理论，而不是先写代码再补一套 rubric。

\subsection{什么时候评：训练、开发、模型选择与生产}

评估时点会改变数据和指标。训练期需要便宜、频繁、能定位问题的 regression set；开发期比较 prompt、tool policy、memory 与 harness；模型选择要在同协议下做公平比较；部署后则监控真实分布、延迟、成本、安全和 drift。课堂 slide 5 把这些阶段画在同一生命周期上，提醒团队不要把 benchmark 当成发布前一次考试。

\begin{table}[H]
\centering
\begin{tabular}{p{2.2cm}p{3.5cm}p{4.2cm}p{2.5cm}}
\toprule
\textbf{阶段} & \textbf{核心问题} & \textbf{推荐证据} & \textbf{典型频率} \\
\midrule
训练 & 更新是否学到目标能力 & 自动 verifier、分桶 reward、轨迹统计 & 每个 checkpoint \\
开发 & 组件改动是否真实改善 & 固定模型、harness swap、错误归因 & 每次变更 \\
选择 & 哪个组合更适合目标场景 & 盲测、置信区间、成本/质量前沿 & 发布候选 \\
部署 & 线上是否持续满足边界 & 真实任务、漂移、安全、回滚指标 & 持续监控 \\
\bottomrule
\end{tabular}
\caption{不同生命周期阶段的 evaluation 任务}
\end{table}

\begin{warningbox}{同一测试集不能承担所有职责}
频繁调参会让开发者对固定测试集过拟合；生产流量又会出现训练时没有的工具版本、权限和用户行为。应把 smoke/regression、研究 holdout、上线 canary 和事故回放分层管理，并记录数据版本。一个分数在训练期有用，不代表足以批准高风险部署。
\end{warningbox}

\subsection{本章小结}

Agent evaluation 的对象是“模型 + harness + tools + environment”的运行系统。它贯穿全生命周期，并必须为具体决策服务。接下来先建立任务与评分类型，再讨论 Agent-specific taxonomy 和 outcome validity。

